\documentclass[10pt,a4paper]{amsart}
\usepackage[margin=19mm]{geometry}
\usepackage{amsmath,amssymb,mathtools}
\usepackage[scr=rsfs,cal=euler]{mathalfa}
\usepackage{xfrac}
\usepackage[unicode,hidelinks,bookmarks=false]{hyperref}
\newcommand{\ie}{i.e.\ }
\newcommand{\cf}{cf.\ }
\newcommand{\resp}{respectively\ }
\newcommand{\Subsection}{\subsection}
\providecommand{\nobreakdash}{\nobreak\hbox{-}}
\setlength{\emergencystretch}{2em}
\multlinegap0pt
\usepackage{bm}
\newtheorem{lemma}{Lemma}
\title{Three heteroclinic orbits induce a countable family of equivalence classes of regular flows}
\author{Elena Gurevich}
\date{Source: arXiv:2603.07470v1 (CC BY 4.0)}
\begin{document}
\maketitle

\noindent\emph{Source and licence.} Three heteroclinic orbits induce a countable family of equivalence classes of regular flows. Version arXiv:2603.07470v1 (CC BY 4.0), distributed by arXiv under the Creative Commons Attribution 4.0 International licence, http://creativecommons.org/licenses/by/4.0/, as recorded in the arXiv licence metadata for this exact version and shown on the abstract page https://arxiv.org/abs/2603.07470v1. Full licence notice: \texttt{licenses/arxiv-2603.07470-CC-BY-4.0.txt}.\par\smallskip

\noindent\textit{Editorial scope.} Counted unit: the article's lemma dehn (source0.tex lines 320--329). The article's complete original proof of it is delivered with it (source0.tex lines 330--367, one \texttt{proof} environment of 38 lines). No mathematics is written, repaired or re-derived: the statement and the proof are the article's own lines, in the article's own notation, and no retained line is substituted or re-flowed.  No further source-local context is needed: every cross-reference of every retained line resolves inside the two delivered spans.  Nothing was omitted from any retained span, so the package carries no omission record.  No retained line contains a citation, so the wrapper carries no bibliography and no bibliography entry.  No retained line contains a figure environment, an image-inclusion command or drawing code, so no image is delivered and the package declares no asset dependency.  Editorial formatting only, in addition: an AMS \texttt{amsart} preamble that restates the article's own declarations and macros used by the retained lines, namely the environments \texttt{lemma} on separate counters, together with the standard packages those lines need.  The retained environments print in this document as Lemma 1 in order of appearance, whereas the article prints the same items with the numbering of its own full document.  The complete native source of the article as distributed in the arXiv e-print for this exact version is bundled unchanged as \texttt{sources/195945/source0.tex}.
\begin{lemma}\label{dehn}
Let $\mathbb B^2_r=\{x\in \mathbb R^2|\ |x|=r\}$, $0<r_1<r<1$; $x\in \partial \mathbb B^2_r$ --- some point and $g: \mathbb B^2_r\times [0,1]\to \mathbb B^2_r\times [0,1]$ --- a homeomorphism such that $g(\{O\}\times [0,1])=\{O\}\times [0,1]$. Then there exists a homeomorphism $h:\mathbb B^2\times [0,1]\to \mathbb B^2\times [0,1]$ with the following properties:
\begin{enumerate}
\item $h|_{\mathbb B^2_{r_1}\times [0,1]}={\rm id}$;
\item $h(\mathbb B^2_r\times [0,1])=\mathbb B^2_r\times [0,1]$;
\item $h(g(x\times [0,1]))=x\times [0,1]$;
\item $h|_{\partial \mathbb B^2\times [0,1]}={\rm id}$.
\end{enumerate}

\end{lemma}

\begin{proof} Set $l=x\times [0,1], l'=g(l)$. Let $e:[-1,1]\times [0,1]\to \partial \mathbb B^2_r\times [0,1]$ be a topological embedding such that $e(\{0\}\times [0,1])=l$. Set $N_l=e([-1,1]\times [0,1])$ and denote by $e', N_{l'}$ a similar embedding and its image, which is a compact neighborhood of the arc $l'$ in $\partial \mathbb B^2_r\times [0,1]$.
The set $$C_l={\rm cl} (\partial \mathbb B_r\times [0,1]\setminus N_l)$$ is homeomorphic to the disk $\mathbb B^2$. By Alexander Theorem, the restriction of the homeomorphism $e'e^{-1}: N_l\to N_{l'}$ to $\partial N_l$ extends to a homeomorphism $C_l\to C_{l'}$. Therefore, there exists an orientation-preserving homeomorphism
$$\psi: \partial \mathbb B_r\times [0,1]\to \partial \mathbb B^2_r\times [0,1]$$
such that $\psi(l)=l'$. 

Set $$P_+=(\mathbb B^2\setminus {\rm int}\, \mathbb B^2_{r}) \times [0,1],\  P_-=(\mathbb B^2_{r}\setminus {\rm int}\, \mathbb B^2_{r_1}) \times [0,1].$$

 $P_+, P_-$ are homeomorphic to the solid torus   $\mathbb B^2\times \mathbb S^1$. We extend $\psi$ to a homeomorphism $\Psi_\pm:\partial P_\pm\to \partial P_\pm$ that will be identity on $\partial B_r\times [0,1], \partial B_{r_1}\times [0,1]$ and could be extended to a homeomorphism $\Phi_\pm: P_\pm\to P_\pm$. 

Let $\psi_{i}=\psi|_{\partial \mathbb B^2_r\times \{i\}}, i\in \{0,1\}$, and denote by $\psi^t_i: \partial \mathbb B^2_r\times \{i\}\to \partial \mathbb B^2_r\times \{i\}$ an isotopy connecting $\psi^1_i=\psi_i$ with the identity map $\psi^0_i$.  Set $$\tau_+(|x|)=\frac{1-|x|}{1-r},$$ and  define the   homeomorphisms $\Psi_+: \partial P_+\to \partial P_+$ by the relations

$$
\Psi_+(x)=\begin{cases}\psi(x),\ x\in \partial \mathbb B^2_r\times [0,1];\cr
x,\ x\in \partial \mathbb B^2\times [0,1];\cr
\frac{|x|}{r} \psi^{\tau_+(|x|)}_i(\frac{rx}{|x|}),\ x\in \mathbb B^2\setminus \mathbb B^2_r\times \{i\}, i\in \{0,1\}.
\end{cases}
$$

Let  $\rho_{x}$  be a radial  segment of   $\mathbb B^2$ joining the point $x\in \partial \mathbb B^2_r$ with a point  $x'\in \partial \mathbb B^2$. Then the union 

$$\mu=l\cup \{\rho_x\times \{0,1\}\}\cup \{x'\times [0,1]\}$$

is the meridian of solid torus $P_+$. A curve $\Psi_+(\mu)\subset \partial P_+$ is a closed curve that intersects a canonical longitude $\partial \mathbb B^2\times \{0\}$ of the torus $\partial P_+$ at exactly one point. Then there  are two possibilities. 

1)   $\Psi_+(\mu)$ is a meridian of $P_+$, in this case there exists a homeomorphism $\Phi_+: P_+\to P'_+$ such that $\Phi_+|_{\partial P_+}=\Psi_{+}|_{\partial P_+}$.

2) There is a Dehn twist $\tau: \partial P_+\to \partial P_+$ with a support $(\mathbb B^2\setminus {\rm int}\, \mathbb B^2_r)\times \{0\}$ such that
$\tau(\Psi_+(\mu))$ is  a meridian of $P_+$.  In this case the composition $\tau \Psi_+$ sends $l$ to $l'$ and  extends to the desired homeomorphism $\Phi_+: P_+\to P'_+$.

 Similarly, one may construct a homeomorphism $\Phi_-: P_-\to P_-$ such that $\Phi_-|_{\partial P_+\cap \partial P_-}=\Phi_+|_{\partial P_+\cap \partial P_-}=\psi_{\partial P_+\cap \partial P_-}$. Then the  desired homeomorphism $h$ is defined by the formula

\begin{equation}
h(x)=\begin{cases}\Phi_+(x), x\in (\mathbb B^2\setminus {\rm int}\, \mathbb B^2_{r}) \times [0,1];\cr
\Phi_-(x), x\in (\mathbb B^2_r\setminus {\rm int}\, \mathbb B^2_{r_1}) \times [0,1];\cr
x,\ x\in \mathbb B_{r_1}\times [0,1].\end{cases}
\end{equation}

\end{proof}

\end{document}
